\section{Introduction to the Topic and Historical Context}

\subsection{Introduction}

It is the autumn of 1975. Over the past decades, the ideological map of Latin America had evolved into one of the most contested battlegrounds of the Cold War. By definition, the Cold War constitutes a geopolitical, ideological and economic rivalry between the United States and the Soviet Union, and is specifically marked by tensions, proxy wars, and an arms race without direct large-scale open conflicts between the superpowers.

On the one hand, in the 1960s, elected or semi-democratic governments across the Southern Cone have fallen into the hands of military groups, such as Brazil in 1964, Bolivia in 1971, Uruguay and Chile in 1973. Following that, each new junta (military or political group that rules a country after taking power by force) has drawn closer to Washington and to one another. Moreover, Argentina remains, for now, under the fragile civilian government of Isabel Perón, but political violence is escalating rapidly.

On the other hand, intelligence gathered by Soviet and Cuban services indicates that the region's military governments are no longer acting only within their own borders:

\begin{itemize}
  \item security officials from different countries are exchanging files on exiled dissidents,
  \item tracking political refugees across frontiers,
  \item and in some cases, assisting one another in surveillance and harassment of political opponents.
\end{itemize}

What has so far been an informal, bilateral practice of repression is beginning to show signs of a shift. Intelligence services across the Southern Cone have started exchanging information on wanted individuals with increasing regularity, and preliminary contacts suggest an appetite for something more permanent: shared watch-lists, surveillance of exiles crossing borders, and informal agreements to detain or hand over targets on request. Whether this remains ad hoc or hardens into a standing, coordinated apparatus is a matter this cabinet is now in a position to shape.

Nothing about that outcome is fixed yet, it is precisely what this committee, alongside the three others in this crisis, exists to decide.

At its core, this crisis is a struggle for power. Indeed, Latin America's military leaders are moving to secure their hold on government and eradicate every trace of organized left-wing opposition, while revolutionary movements across the continent are moving in the opposite direction: planning to overthrow those same governments and seize control of the region for themselves. Circling both sides are two intelligence services bound to neither; the American CIA and the Soviet KGB,  each pursuing its own interests in Latin America, free to work with, around, or against anyone the moment it serves that interest. This is a crisis of military combatants, political mobilization, treason, and intelligence sharing, and it will not be decided by any single battle. It will be decided by the accumulated weight of the alliances and decisions each cabinet makes as the crisis unfolds, and by the confrontations that break out \emph{within} each cabinet, which will make cooperation, even among supposed allies, far harder than delegates might expect.

\subsection{Key terms}

\textbf{KGB (Komitet Gosudarstvennoy Bezopasnosti / Committee for State Security):} The KGB is the Soviet Union’s principal security agency, established on March 13, 1954. It merges foreign intelligence, domestic secret police work, counterintelligence, and the Communist Party’s political enforcement arm under one roof. Operating from its infamous Lubyanka headquarters in Moscow, it plays a decisive role in shaping Soviet policy toward Latin America, at times even taking precedence over the Ministry of Foreign Affairs itself.

\textbf{DINA (Dirección de Inteligencia Nacional):} DINA is Chile’s secret police under Augusto Pinochet, established in the wake of the 1973 coup. It handles the surveillance, detention, and elimination of political opponents, including MIR leaders such as Dagoberto Pérez Vargas. DINA stands out as a textbook example of the national security apparatuses whose activities, and potential cross-border coordination, lie at the center of this crisis.

\textbf{Departamento América (DA):} Departamento América is a division within the Cuban Communist Party’s Central Committee, tasked with coordinating Cuba’s support for revolutionary and guerrilla movements across Latin America. It provides training, funding, and ideological backing, effectively turning Havana into a revolutionary nerve center for the region, linked to Soviet strategy, yet operating on its own terms.

\textbf{Unidad Popular (Popular Unity):} Unidad Popular, or Popular Unity, is the coalition of socialist, communist, and other left-wing parties that brought Salvador Allende to power in Chile through democratic elections in 1970. It marked the first time a self-declared Marxist government had come to power via the ballot box rather than armed revolution. Its violent overthrow in the 1973 coup remains a foundational event, shaping the ideological stakes of this entire crisis.

\textbf{Détente:} A French term meaning "relaxation," détente describes the general easing of tensions between the United States and the Soviet Union during the 1970s. Even with agreements like SALT I and the 1975 Helsinki Accords, however, it coexists uneasily with continued ideological and covert rivalry in places like Latin America. Both superpowers continue jockeying for influence abroad, but carefully avoid any direct military clash.

\textbf{Junta:} A junta is a military or political group that seizes control of government, usually by forcibly overthrowing the existing regime. The term lies at the heart of this crisis: during the 1970s, several Latin American countries, Chile, Brazil, Uruguay, Bolivia, and Argentina; were either under junta rule or dangerously close to it.

\textbf{Active Measures:} Active Measures are a KGB specialty: covert political influence operations aimed at discrediting rival powers, particularly the United States, without triggering a direct military engagement. These include disinformation campaigns, forged documents, propaganda routed through front organizations, and deliberate manipulation of foreign public opinion.

\textbf{Guerra Sucia (Dirty War):} The Guerra Sucia, or "Dirty War," refers to the wave of state-sponsored violence, forced disappearances, and repression carried out by Argentina’s military and allied paramilitary groups, such as the Argentine Anticommunist Alliance (AAA), against real or suspected left-wing dissidents.

\textbf{Domino Theory:} The Domino Theory is a Cold War-era geopolitical assumption, largely associated with U.S. foreign policy, which holds that if one state fell to communism, neighboring states would topple in a chain reaction. This conviction provides the rationale for Washington’s backing of anti-communist governments and juntas throughout Latin America as well as Southeast Asia.

\textbf{Plausible deniability:} Plausible deniability is an intelligence and political practice in which a government or agency structures its covert operations so that its involvement cannot be conclusively proven or officially acknowledged if things go wrong.

\subsection{Historical context}

The Russian acronym KGB actually refers to the “\emph{Komiteth Gosudarstvennoy Bezopasnosti}”, and is also defined as the Committee for State Security. It was founded on March, 13th 1954, under the Union of Soviet Socialist Republics’ (USSR) Council of Ministers’ authority. Historically, the KGB creation resulted both in the aftermath of Stalin’s death and its predecessor’s dismantling, the \emph{MGB (Ministerstvo Gossoudarstvenno}ï \emph{Bezopasnosti)}. From its headquarters located at the Lubyanka in Moscow, it operated simultaneously as a:

\begin{itemize}
  \item Foreign intelligence service,
  \item Domestic secret police,
  \item Counterintelligence agency,
  \item Political enforcement arm of the Soviet Union Communist Party
\end{itemize}

By 1975, under Yuri Vladimirovich Andropov, longest-serving Chairman, appointed in May 1967,  the KGB had been transformed from a blunt instrument of Stalinist terror into a sophisticated engine of ideological warfare, covert influence, and strategic deception, operating parallel to, and often ahead of, the Soviet Foreign Ministry itself.

\subsubsection{The Soviet Intelligence system and the extension of its influence across Latin America}

Some of what this cabinet decides may not remain hidden indefinitely. Concerns have occasionally been raised, within intelligence services generally, about the long-term security of internal records, and it is plausible that materials documenting this period could surface at some later point, through channels not yet known.

Moreover, Soviet intervention in Latin America is the result of twenty years of influence. It was actually the KGB, rather than the Foreign Ministry of the Soviet Union, that played the leading role in forming Moscow's policies towards Latin American countries (Siegel, 2005). Second, Cuban alliance with Moscow after 1959 is, in the KGB assessment, a strategic victory. Indeed, the coming to power of Fidel Castro, a reliable pro-Soviet ally, was of great help in enhancing the prestige and image of the revolutionary Soviet Union in the Third World. It is the experience that Andropov, Kryuchkov, and Leonov bring to this crisis situation: Latin America is not on the periphery of Soviet policy but rather one of the few regions where the USSR has managed to claim ideological advantage over Washington.

Following that, the KGB’s assets were therefore never limited to espionage in the conventional way. Effectively, the cabinet’s disinformation directorate, ran sustained “active measures” campaigns designed to discredit Western (and notably American) influence in the region. For example, fabricated atrocity stories and well rounded propaganda placed throughout front organizations (Federation of American Scientist, 1997). Thus, these are not abstract Cold War relics; but active measures capacity is one of the concrete tools this KGB office can deploy, and its use elsewhere in the region. To illustrate, including disinformation efforts built around Nazi émigrés like \emph{Walther Rauff} can highlight how far Moscow has been willing to blur the line between intelligence work and political warfare (Orton, 2025).

\subsubsection{Chile: Salvador Allende’s Popular Unity government}

Salvador Allende’s Popular Unity Government, which was elected in 1970, represented the first time a self-declared Marxist had won power in Latin America, through the conventional way of the ballot bow rather than with a revolution. In parallel, a viability assessment of the Popular Unity project, commissioned by Yuri Andropov himself and known as the Informe Andropov, concluded that Chile's economic demands exceeded what the USSR could sustainably provide, showing that Moscow’s support was never unconditional.  This tension of real solidarity constrained by real Soviet caution, constitutes the conflict that will have to be managed going forward. Thus, it is precisely the reason why Mikhail Suslov and Boris Ponomarev’s ideological pressures for a “more decisive commitment” to endangered allies, matters as a political force.

However, all of these efforts were ineffective. The Chilean military, under the leadership of General Augusto Pinochet, toppled Allende on September 11, 1973 in a coup which had been strongly supported by the United States. In fact, historians working with the National Security Archive have documented that President Nixon ordered the CIA in 1970 to "make the economy scream" in order to block Allende's inauguration in the first place (National Security Archive, 2022). The more specific position of Henry Kissinger as the US Secretary of State was chairing the Washington Special Actions Group that helped to coordinate U.S. aid for the newly formed junta in less than a day after the military coup.

\subsubsection{Cuba, the Guerilla networks and the limits of Soviet control}

Fidel Castro serves as both head of state and head of government of Cuba, a position from which he has positioned the island as one of the most active defenders of revolutionary Marxism in the hemisphere, through direct material support to allied movements rather than rhetoric alone. Castro maintains formal relations with the Soviets, while relations with several neighboring governments range from cautious to openly hostile, depending on each government's own tensions with leftist movements. This dual role, engaging in ordinary statecraft while sponsoring armed revolutionary struggle abroad, creates ongoing friction with neighboring states and places Havana at the center of regional tensions that must be navigated.

Departamento América (DA), an arm of the Cuban Communist Party's Central Committee, coordinated support to revolutionary and guerrilla organizations across the hemisphere. It is a central conduit for Cuba's training, funding, and ideological support to allied movements (Hudson, 1988), and a dense network of Cuban-linked organizations extending from Central America to the Southern Cone (Masterson, 1997).

This activity connects directly to the wider Soviet effort: intelligence services view Cuba as a force-multiplier capable of moving people, funds, and material in ways less directly available to others, often with a degree of deniability that is highly valued. This dynamic ties closely to parallel discussions on hemispheric strategy, where decisions on pace and risk tolerance will constrain or embolden what Havana is willing to attempt independently.

The Chilean MIR (Movimiento de Izquierda Revolucionaria) and the Argentinian ERP (Ejército Revolucionario del Pueblo) are the two guerrilla organizations explicitly at stake. Following the change of government in Chile, the MIR began cautiously accepting logistical assistance from outside Cuba.

The region's revolutionary left had also been organizing itself independently. The joining of the Argentine ERP, the Chilean MIR, the Bolivian ELN, and the Uruguayan Tupamaros produced an attempt at a regional revolutionary alliance operating independently of Moscow's approval.

Cuba's willingness to support such organizations independently of established timelines and risk tolerance remains the main source of friction. Manuel Piñeiro, Fidel Castro, and Raúl Castro consistently push for more direct action than more cautious, orthodox voices within the Soviet group are prepared to authorize.

\subsubsection{Coordinated repression}

Although repression has remained primarily a national undertaking, increasing evidence points to growing cooperation between the security and intelligence services of several Southern Cone governments. Bilateral exchanges of intelligence, the monitoring of political exiles, and informal coordination in the surveillance of suspected subversives have become more frequent as revolutionary organizations operate across national borders. While the extent of this cooperation remains difficult to determine, regional security concerns have created incentives for closer collaboration among anti-communist regimes.

Throughout much of the Southern Cone, military governments have adopted an expansive definition of the “internal enemy”. Beyond armed guerrilla organizations, this concept often encompasses trade unionists, student leaders, political exiles, and individuals perceived to sympathize with Marxist ideology, regardless of whether they have participated in armed struggle. As opposition movements increasingly transcend national frontiers, intelligence-sharing and cross-border coordination have emerged as practical tools for governments seeking to prevent suspected opponents from exploiting jurisdictional boundaries.

These developments unfold within the broader geopolitical context of the Cold War. The Soviet Union and Cuba maintain political, diplomatic, and intelligence interests throughout Latin America, while the United States continues to view regional security through the lens of communist containment.

\subsubsection{Regional Tensions}

By November, Latin America has become one of the tensest areas within the Cold War. Through multiple coups that had replaced elected or constitutional governments, and military regimes which enacted the repression of socialist movements, the continent serves as one of the biggest hot spots in the ideological clash between communism and capitalism. With Brazil having fallen under military rule over a decade ago, Bolivia, Uruguay, and Chile soon followed.

For the Soviet Union, these developments represent more than isolated domestic crises. The Southern Cone progressively appears to be evolving into a region dominated by governments willing to cooperate against socialist organizations. At the same time, revolutionary organizations are attempting to overcome national divisions of their own. Guerrilla movements operating in Uruguay, Bolivia, Argentina, Chile, and elsewhere increasingly recognize that isolated insurgencies are unlikely to survive against coordinated military governments. Consequently, Latin America has entered a period in which both governments and revolutionaries are pursuing regional coordination, with each of them seeking to outpace the other.

This growing competition forms the central fault of this crisis. On one side, military governments seek to consolidate anti-communist rule and possibly eliminate any organized opposition. On the other side stands revolutionary organizations, attempting to preserve communism through insurgency, underground movements, and international solidarity. The Soviet KGB and the American CIA each have chosen their sides, finding allies in Latin America whose objectives frequently overlap with their own. While publicly both superpowers’ relationship seem to be in a calmer period, these tensions can easily threaten the period of détente that they are currently in.

\minorheading{Tensions with the United States}

The struggle in the continent cannot be separated from the broader rivalry between the United States and the Soviet Union. Since the Cuban Revolution, the White House has gradually viewed Latin America as a possible starting point for the Domino Theory to take effect. Due to this belief, Washington has a history of supporting governments that promised stability and anti-communist legislation. This was made clear with the U.S.’ backing of the Khmer Republic in Asia, and their decision to quickly establish formal diplomatic relations and provide economic and military aid to Pinochet's junta. Meanwhile, Moscow regarded Latin America as a key area where socialist influence still possessed significant opportunities for expansion. Yet, by November of 1975, neither superpower has engaged in direct military confrontation.

Both nations remained formally committed to the policy of détente, engaging in strategic dialogue while secretly competing for influence in multiple regions around the globe. While both nations seemed to be working on improving their relations through treaties such as the Helsinki Accords, none were willing to give up their ambitions in influencing other nations. This contradiction between public sentiment and state level orders and priorities meant that conflict increasingly occurred through intelligence operations and economic assistance, rather than conventional warfare.

Growing distrust directly threatens to undermine this fragile relationship, as in the United States, there is growing dissent against détente. Figures like Ronald Reagan, the Governor of California, and Senator Henry M. "Scoop" Jackson represented a growing frustration with détente which ultimately grew after the American loss in Vietnam. Soviet policymakers have similar worries, with many believing that American interventions against socialist governments and their support for anti-communist regimes revealed a lack of commitment to improving relations.

\minorheading{Détente and the International Situation}

By November 1975, the Soviet Union and the USA have found themselves in one of the most contradictory moments of the Cold War. In August, 35 states had signed the Helsinki Final Act (also known as the Helsinki Accords) to formally recognize the post-war European order while promoting cooperation in security, economics, and human rights. To many observers, the agreement represents the peak of détente, as the Accords encompassed the importance of diplomacy and cooperation.

Beneath the appearance of stability, tensions continued to accumulate. From August onwards, different negotiations have begun to stall. Notably, talks regarding the second Strategic Arms Limitation Treaty (SALT II) remain ongoing but increasingly tense and difficult. Simultaneously, crises in Africa, the Middle East, and Latin America have demonstrated that ideological competition has merely shifted to another playing stage rather than disappearing altogether.

\minorheading{The Rise and Fall of Latin American Governments}

These countries have strong socialist and military movements, yet each of them find themselves in different situations.

\textbf{Paraguay:} Since 1954, the country has been under the dictatorship of Alfredo Stroessner. As the longest-standing dictatorship in the Southern Cone, it is the most politically and economically stable compared to its neighbours. With extensive internal security and intelligence services, Paraguay increasingly serves as a meeting point for other military regimes. For both revolutionary organisations and foreign intelligence services, Paraguay represents one of the most hostile operating environments in Latin America.

\textbf{Bolivia:} Stemming from the 1971 coup, Bolivia has been governed by a military regime which is firmly committed against communism. Although the defeat of Che Guevara’s guerilla campaign had weakened organized military movements in the country, left-wing resistance hasn’t disappeared. Organizations such as the Ejército de Liberación Nacional (ELN) continue to seek opportunities to reorganize the leftist movement. Bolivia, therefore, remains unstable, with its government balancing internal repression with persistent social unrest.

\textbf{Brazil:} Similarly to Paraguay, Brazil has been under an anti-communist dictatorship since 1964. Its intelligence and security services possess considerable experience in surveillance, censorship, and dealing with counter-insurgency. While large-scale guerilla activity has been largely suppressed by 1975, authorities continue to cooperate closely with regional partners in monitoring exiles.

\textbf{Peru:} Compared to other nations, Peru’s balance of regional power appears more ambiguous. When Francisco Morales Bermúdez came into power in August, he initially portrayed himself as independent. However, he has progressively aligned with neighboring right-wing, anti-communist military dictatorships in South America (such as Chile, Argentina, and Brazil). However, due to the sudden changes in economic policies which affect food prices and labourers, there is growing social unrest.

\textbf{Argentina:} Out of all the nations, Argentina may represent the greatest uncertainty in confronting the region. With Isabel Perón coming into power the previous year, the country finds itself amid regular clashes between left-wing guerrilla groups and right-wing death squads, and great economic instability. Still amidst the Guerra Sucia, left-wing guerilla groups such as the ERP and Montoneros found themselves organizing armed campaigns against the government and funded the Argentine Anticommunist Alliance (AAA). Many observers believe that military intervention within the country has become a matter of when rather than if, making Argentina a possible principal battleground if the US or Military regime alliance decide to directly intervene.

\textbf{Ecuador:} Ecuador exists as a unique case. It remains relatively removed from escalating confrontation between military governments and revolutionary organizations. The country, however, is entering a phase of high instability, with the aftermath of the September coup leaving the government in a fragile position, while the mobilization of trade unions and workers are beginning to increase.

\textbf{Chile:} Two years after the overthrowing of Salvador Allende’s government, Chile has become the symbol of military anti-communism in Latin America. Pinochet’s regime has dismantled and cracked down on multiple left-wing political organizations, having imprisoned or exiled thousands of opponents. Some of these include key MIR leader Nelson Gutiérrez, who was given the Vatican's diplomatic residence in Santiago, and the killing of another MIR leader Dagoberto Pérez Vargas during a raid conducted by the DINA (secret police). For revolutionary organizations, Chile demonstrates the consequences of military consolidation, while offering a model on how to conduct political repression for neighboring governments.

\textbf{Uruguay:} Once regarded as one of the continent’s most stable democracies, the 1973 coup marked the beginning of a dramatic transformation within the country. The military regime has successfully weakened the Tupamaros guerilla movement through widespread arrests, intelligence operations, and emergency security legislation. Despite this success, the country suffers political instability. While there was no open street chaos, political repression continues, while many Uruguayan exiles remain active abroad.
